\subsection{Experiment details} \label{sec:experiment_details}

We outline two "studies" - one focusing on the composition of transformations, the \textbf{Compositional Generalization} \textit{(CompGen)} study, and the other on environment variations (with fixed transformation sequence): the \textbf{Environmental Generalization} \textit{(EnvGen)} study. In each case, we measure the capacity of models to perform tasks that differ in some way from the tasks seen during training---either with respect to transformation composition in the \textit{CompGen} study or with respect to environment parametrization in the \textit{EnvGen} study.

For both studies, we define different \textit{experiment settings} which focus on different aspects of compositional and systematic generalization. For each \textit{experiment setting} we create 5 "experiments": they are different instances of the \textit{experiment settings} in which we only vary the transformation sequences. This is to ensure that the results we report are robust across different transformation combinations (e.g., translation-based transformation sequences might be easier than rotation-based ones; see Appendix \ref{appendix:transformation_learnability}).

For each experiment, we train on 100,000 unique training samples, and test on two unique and distinct sets of 1,000 samples; one test-set follows the same distribution (i.e., in-domain (ID) set) as the training set, and the other is out-of-distribution (OOD), from which we evaluate generalization.  

In the \textit{CompGen} study, we train and test a model on multiple transformation sequences within a single experiment. This contrasts with the \textit{EnvGen} study, where models are trained on only a single transformation per experiment. To handle this complexity, we provide context that indicates to the model which sequence to perform. We do so by appending a task embedding to the input sequence, similar to the approach in CVR \cite{cvr2022}. This embedding is a sequence of tokens that specifies the transformations in the correct order. For example, the task “translate\_up-rotate\_90” would have an embedding like ['T', 'R'], while the inverse task "rotate\_90-translate\_up" would be ['R', 'T'].
For OOD testing, the models may be evaluated on new or longer sequences of these transformations. However, every individual token would have been seen during training. 

We provide below an outline of each experiment setting for each of our two studies, and refer the reader to the Appendix \ref{appendix:experiment_details} for more details.

\newpage
\paragraph{\textit{CompGen} – Compositional Generalization}
    \begin{itemize}
        \item \textit{C1 – Atomic + Composite → Unseen Composite}: Train on atomic (depth 1) and composite (depth 2) tasks; test on unseen depth 2 composites built out of the same atomic transformations as training.
        \item \textit{C2 – Restricted Composite → Unseen Composite}: Train only on a subset of composites (depth 2); test on unseen depth 2 composites built out of the same atomic as training.
        \item \textit{C3 – Atomic + Composite → Deeper Composite}: Train on atomic and depth 2 composites; test on depth 3 composites built out of the same atomics as training.
    \end{itemize}

\paragraph{\textit{EnvGen} – Environmental Generalization}
\begin{itemize}
    \item \textit{G1 – More Objects}: Train with 1–2 objects; test with 3–4 objects (fixed grid size and object complexity).
    \item \textit{G2 – Larger Grids}: Train on 10×10–15×15 grids; test on 16×16–20×20 grids (fixed object number and complexity).
    \item \textit{G3 – Larger Objects}: Train on objects sized 1×1–5×5; test on objects sized 6×6–10×10 (fixed grid size and number of objects).
    \item \textit{G4 – More Complex Objects}: Train on symmetric, single-colored objects; test on asymmetric, multi-colored (fixed grid size and number of objects).
    \item \textit{G5 – Combined}: Train on simple cases (1–2 objects, 10×10–15×15 grids, symmetric single-colored 1×1–5×5 objects); test on all harder variants simultaneously (3-4 objects, 16x16-20x20 grids, multi-colored asymmetric objects).
\end{itemize}


% \subsubsection{CompGen: Compositionality Generalization}\label{experiments:CompGen}

% For the compositionality study, we perform experiments within the following settings: 

% \begin{itemize}
%     \item \textbf{C1 – Atomic + Composite → Unseen Composite}: Train on atomic (depth 1) and composite (depth 2) tasks; test on novel depth 2 composites. Tests generalization to new compositions.
%     \item \textbf{C2 – Restricted Composite → Unseen Composite}: Train only on a subset of depth 2 composites; test on novel depth 2 composites. Tests generalization from limited composite examples.
%     \item \textbf{C3 – Atomic + Composite → Deeper Composite}: Train on atomic and depth 2 composites; test on depth 3 composites. Tests generalization to greater compositional depth.
% \end{itemize}

% \subsubsection{EnvGen: Environmental Generalization}\label{experiments:EnvGen}

% For the generalization study, we perform experiments within the following settings: 

% \begin{itemize}
%     \item \textbf{G1 – More Objects}: Train on 15×15 grids with 1–2 objects; test with 3–4 objects.
%     \item \textbf{G2 – Larger Grids}: Train on 10×10–15×15 grids; test on 16×16–20×20, with 2 objects.
%     \item \textbf{G3 – Larger Objects}: Train on objects sized 1×1–5×5; test on 6×6–10×10.
%     \item \textbf{G4 – More Complex Objects}: Train on symmetric, single-colored objects; test on asymmetric, multi-colored (2 objects, 15×15 grid, <6×6 size).
%     \item \textbf{G5 – Combined}: Train on simple cases (1–2 objects, 10×10–15×15 grids, symmetric single-colored 1×1–5×5 objects); test on all harder variants simultaneously.
% \end{itemize}



% For the compositionality study, we perform experiments within the following settings: 

% \begin{itemize}
%     \item \textbf{C1 - From Atomic and Composite Tasks to Unseen Composite}: Models are trained on a \textit{mix of individual atomic tasks (depth = 1) and composite tasks (depth = 2)}. The OOD test set consists of unseen composite tasks (depth = 2) composed of the same atomic transformations in novel orders. This setting evaluates whether models can generalize to new compositions when trained on both atomic and composite examples.
%     \item \textbf{C2 - From Restricted Composite Tasks to Unseen Composite}: Models are trained \textit{solely on a restricted set of composite tasks (depth = 2)}. The OOD test set consists of unseen composite tasks (depth = 2) composed of the same atomic transformations in novel orders. This setting evaluates whether models can generalize to new compositions when trained only on other composite examples.
%     \item \textbf{C3 - From Composite Tasks to Deeper Composite Tasks}: Models are trained on \textit{mix of individual atomic tasks (depth = 1) and composite tasks (depth = 2)}. The OOD test set consists of deeper composite tasks (depth = 3) composed of the same atomic transformations. This setting evaluates whether models can generalize to compositions of greater depth than those seen during training.
% \end{itemize}

% \subsubsection{EnvGen: Environmental Generalization Study}\label{experiments:EnvGen}

% For the generalization study, we perform experiments within the following settings: 

% \begin{itemize}
%     \item \textbf{G1 - Generalization to More Objects}: Models are trained on 15x15 grids containing \textit{1 or 2 objects}. The OOD test set consists of grids of the same size containing \textit{3 or 4 objects}. This setting evaluates whether models can generalize to a greater number of objects.

%     \item \textbf{G2 - Generalization to Larger Grids}: Models are trained on grids with sizes ranging from \textit{10x10 to 15x15}. The OOD test set consists of grids with larger sizes, from \textit{16x16 to 20x20}. This setting evaluates whether models can generalize to larger spatial environments while the number of objects is held constant at 2.

%     \item \textbf{G3 - Generalization to Larger Objects}: Models are trained on grids containing objects with dimensions ranging from \textit{1x1 to 5x5}. The OOD test set consists of grids containing objects with larger dimensions, from \textit{6x6 to 10x10}. This setting evaluates whether models can generalize to objects of a significantly different scale.

%     \item \textbf{G4 - Generalization to More Complex Objects}: Models are trained on grids containing \textit{symmetric and single-colored objects}. The OOD test set consists of grids containing \textit{asymmetric and multi-colored objects}. For both sets, all other parameters (2 objects, 15x15 grid size, object dimensions < 6x6) are fixed. This setting evaluates whether models can generalize from simple object features to more complex ones.

%     \item \textbf{G5 - Combined Generalization}: Models are trained on grids with a combination of simple properties: \textit{1-2 objects}, grid sizes from \textit{10x10-15x15}, and objects that are \textit{symmetric, single-colored, and 1x1-5x5}. The OOD test set consists of grids where all properties are more complex: \textit{3-4 objects}, grid sizes from \textit{16x16-20x20}, and objects that are \textit{asymmetric, multi-colored, and 6x6-10x10}. This setting evaluates robust generalization across multiple difficulty axes simultaneously.
% \end{itemize}



